% ======================================================================
%  Manuscript prepared for IEEE Transactions on Smart Grid.
%  Format follows the reference manuscript (IEEEtran, journal, 10 pt,
%  two-column): L. J. Lepolesa et al., arXiv:2511.06398.
%
%  HOW TO ADD THE REMAINING COLBAB GRAPHS (figures 3-9):
%   The system architecture (Fig. 1) and pipeline (Fig. 2) are drawn with
%   TikZ and render automatically; they can be replaced with higher-res
%   PNGs later if desired.
%   1. Create a folder  figs/  next to this file.
%   2. Export each remaining plot from Colab as .PNG and save it with the
%      exact filename shown in the boxed placeholder.
%   3. Recompile - the placeholder is replaced automatically, so the
%      document always compiles even before images are inserted.
%
%  Compile: pdflatex -> pdflatex  (references are hand-written; no .bib).
%  On Overleaf the IEEEtran class is available automatically.
% ======================================================================
\documentclass[10pt,journal]{IEEEtran}

\usepackage{amsmath,amssymb,amsfonts}
\usepackage{graphicx}
\usepackage{textcomp}
\usepackage[utf8]{inputenc}
\usepackage{booktabs}
\usepackage{array}
\usepackage{algorithm}
\usepackage{algpseudocode}
\usepackage{tikz}
\usetikzlibrary{arrows.meta}

\newtheorem{proposition}{Proposition}[section]

\graphicspath{{figs/}}

% ---- Conditional figure include: shows a boxed placeholder until the PNG exists
\newcommand{\paperfig}[4][\columnwidth]{%
  \begin{figure}[!t]
    \centering
    \IfFileExists{figs/#2}{%
      \includegraphics[width=#1]{#2}%
    }{%
      \fbox{\parbox[c][0.45\linewidth][c]{0.92#1}{%
        \centering\small\itshape INSERT COLBAB GRAPH HERE\\[3pt]
        \upshape\footnotesize save as \texttt{figs/#2}}}}%
    \caption{#3}\label{#4}
  \end{figure}}

\newcommand{\D}{\Delta}
\newcommand{\util}{\mathrm{util}}
\newcommand{\norm}[1]{\mathrm{norm}\{#1\}}

\begin{document}

\title{Generalized $N$-Network EV Charging Pricing with Distance-Aware
Collaborative DRL and Renewable Net-Load Balancing}

% AUTHOR BLOCK - clean two-author header, matching the reference manuscript.
% Re-add affiliations/e-mails at submission time, e.g.:
%   \thanks{Srilekha~R. and Pavithra~S. are with the Department,
%   Institution, City, Country (e-mail: \{author1, author2\}@institution.edu).}
\author{Srilekha~R.,~and~Pavithra~S.}

\maketitle

% ======================================================================
\begin{abstract}
Collaborative deep deterministic policy gradient (DDPG) agents have been
shown to balance EV charging load across two adjacent distribution
networks. That formulation is restricted to exactly two networks, assumes
drivers are equally willing to charge at any neighboring network
regardless of distance, and neglects on-site renewables. We generalize the
balancing signal to an arbitrary number of networks (with an exact
reduction to the pairwise formula when $N{=}2$, proven algebraically and
verified numerically at tolerance $10^{-7}$), introduce a
Haversine-distance exponential decay that weights each neighbor's influence
by physical proximity, and reformulate pricing and demand response in terms
of a solar/wind-adjusted net load with a $5\%$ baseload floor. On a real
Tetouan, Morocco-grounded four-network testbed, the combined framework
reduces the maximum--minimum cross-network utilization imbalance by
$19.03\%$ relative to uncoordinated pricing and by $15.28\%$ relative to the
original pairwise formulation, with the ranking preserved across four
seasons and independent DDPG training initializations. A four-way ablation
attributes most of the gain to renewable net-load integration. A re-analysis
of the base paper's own two-network data shows that a structural EV
peak-uniformity limitation is inherited from its independent min-max
normalization, not introduced by our extensions.
\end{abstract}

\begin{IEEEkeywords}
Electric vehicle charging, dynamic pricing, collaborative deep
reinforcement learning, DDPG, multi-network load balancing, distance decay,
renewable net load, smart grid.
\end{IEEEkeywords}

% ======================================================================
\section{Introduction}\label{sec:intro}
\IEEEPARstart{T}{he} electrification of road transport is progressing
rapidly worldwide, and with it the concentration of a new, large, and
time-flexible electrical load on urban distribution systems. Global sales
of battery-electric vehicles grew by more than half in 2022, crossing ten
million units in a single year for the first time, a trajectory summarized
in the international outlook reported alongside the base
framework~\cite{lepolesa2026}. Because most charging happens at home and in
workplaces at the end of the working day,
EV charging demand coincides with the existing evening peak in many
feeders, elevating peak-to-average ratios, accelerating transformer and
feeder aging, increasing line losses, and in the limit forcing costly
network reinforcement that would otherwise be unnecessary
\cite{clement2010,richardson2012tsg,richardson2012tpwrs,oconnell2014}. The
magnitude of this stress depends fundamentally on how charging is
coordinated: uncontrolled charging has long been shown to degrade voltage
profiles and overload equipment
\cite{clement2010,richardson2012tsg}, whereas controlled and price-guided
charging can keep the same feeder within statutory limits while
accommodating far higher EV penetrations
\cite{richardson2012tsg,richardson2012tpwrs,oconnell2014,gan2013,callaway2011}.

Grid-impact results are typically reported at the feeder level, and for
good reason: distribution transformers and feeders are sized for their
historical load mix, leaving little thermal headroom for an entire
light-duty fleet. Uncoordinated evening charging has long been associated
with accelerated transformer aging, undervoltage, and feeder overloading
on low-voltage networks~\cite{clement2010,richardson2012tsg}, and the
feeder-reconfiguration literature treats load-carrying capacity as the
binding constraint on new load integration~\cite{baran1989}. Transmission
corridors are planned with larger margins, but primary and secondary
feeders rarely double as hosting assets for fleet-scale electrification.
The relevant decision variable in this paper is therefore the
\emph{distribution-level} load shape, and specifically its cross-network
spread, rather than aggregate system demand.

Retail electricity pricing is the most widely deployed instrument for
steering this demand. Yet evidence from real implementations is cautionary.
Large-scale field experiments show that static time-of-use (TOU) tariffs
shift charging into the off-peak window but simultaneously create a
``shadow peak'' at the window boundary, which in some demonstrations
exceeds the original evening peak~\cite{bailey2025,samende2025,jones2022}.
Aggregating fifteen pricing pilots, Faruqui and Sergici report that plain
TOU designs shave peak demand by only $3$--$6\%$, while finer-grained
dynamic tariffs reach $13$--$44\%$ when paired with enabling automation
\cite{faruqui2010}. Borenstein further shows, from a long-run welfare
perspective, that coarse two-part tariffs capture only a small fraction of
the efficiency gains of fast, real-time retail pricing~\cite{borenstein2005}.
Collectively these works establish two requirements for effective EV tariff
design: prices must adapt continuously to grid state, and they must be
allowed to differ across space as well as time.

The practical stakes are large. Distribution-operator deployment studies,
reported in the base framework~\cite{lepolesa2026}, anticipate that a
substantial fraction of primary feeders will require reinforcement by the
mid-2030s--2040s if fleet charging continues to be attracted to the same
evening window. In a solar-rich grid the marginal cost of a midday charge
is close to its opportunity floor, whereas an evening charge effectively
prices in that reinforcement; even a moderate relocation of demand
therefore has material economic value. This is the economic argument for
pricing that varies in \emph{time} and \emph{place} simultaneously--the
two degrees of freedom exploited by the collaborative framework below.

Deep reinforcement learning (DRL) has emerged as the natural framework for
the first requirement. Deep Q-networks~\cite{mnih2015}, and in particular
the continuous-control algorithms deep deterministic policy gradient
(DDPG)~\cite{lillicrap2016}, Soft Actor-Critic~\cite{haarnoja2018}, and
Proximal Policy Optimization~\cite{schulman2017}, have enabled agents to
learn price or dispatch policies directly from raw load and market states.
Applied to EV pricing, DDPG has been used for aggregator tariff
design~\cite{qiu2020tia,liu2021access}, charging-station revenue
maximization~\cite{zhao2022tte}, demand-responsiveness pricing
\cite{lai2022tsg,cui2021tsg}, vehicle-to-grid scheduling~\cite{aljafari2023},
and, more recently, coordination across multiple market or charging
entities~\cite{ye2022,yang2024,paudel2023}. Surveys of reinforcement learning
for demand response summarize this body of work and repeatedly identify two
cautionary patterns: most systems are single-agent and single-network, and
when prices are demand-dependent there is a real risk of merely shifting
rather than shaving peaks~\cite{vazquez2019}.

Taken together, the literature suggests a clean separation of concerns that
this work adopts: gradient-boosted forecasting supplies accurate day-ahead
net-load trajectories~\cite{chen2016,ke2017,hong2016,wang2016recency}; a
closed-form, price-sensitive demand model translates tariff signals into
fleet-level charging decisions~\cite{lepolesa2026}; and DRL agents optimize
the pricing policy that couples the
two~\cite{lillicrap2016,haarnoja2018}. The separation is deliberate: it
keeps every component independently verifiable, and it lets the ablation
study in Section~\ref{sec:results} attribute the measured gains to the
exact stage that produced them.

Among these works, the collaborative framework of Lepolesa \emph{et al.} is
distinct because its explicit objective is not peak-shaving alone but
\emph{load balancing across} distribution networks: two DDPG agents per
network jointly set EV charging tariffs for a residential and a commercial
feeder so that a utilization-difference signal between the two networks is
reduced, in addition to the conventional peak-valley (PV) objective
\cite{lepolesa2026}. In that architecture the DDPG agents act as the pricing
mechanism while the reported utilization metrics are computed from a
closed-form EV demand-response equation; this two-tier evaluation protocol
is faithful to the base method and is reproduced exactly in this work
(Section~\ref{sec:framework}).

Three limitations nonetheless remain, and each is explicitly named as
future work by the original authors~\cite{lepolesa2026}. (i)~The balancing
formulation is restricted to exactly two networks, whereas urban
distribution systems naturally comprise many interconnected zones.
(ii)~EV drivers are assumed to shift charging to any neighbor with equal
willingness, regardless of the physical distance between charging sites--an
assumption that contradicts the distance-sensitive location and mode-choice
behavior established in the spatial-interaction and charging-network
literature~\cite{huff1964,frade2011}. (iii)~The pricing and demand model
considers only conventional load, even though solar and wind penetration
already reshape the net-load envelope of modern
feeders~\cite{hoke2013}.

The spatial dimension is not a cosmetic addition. Charging choices are
jointly driven by price and proximity: drivers measurably discount a
cheaper tariff when it requires a longer detour, a principle formalized by
gravity and Huff-style models of spatial interaction and applied in
charging-network planning~\cite{huff1964,frade2011}. A balancing scheme
that treats every neighbor as equally reachable therefore misdirects
demand, placing excessive weight on distant networks. Replacing that
implicit equal-distance assumption with a calibrated exponential decay over
real zone coordinates is one of the three extensions developed here, and
its quantitative role is isolated in the ablation of
Section~\ref{sec:results}.

This paper closes all three gaps within one consistent framework. The
contributions are as follows.
\begin{enumerate}
  \item \textbf{Generalized $N$-network balancing.} We extend the pairwise
        utilization-difference signal to an arbitrary number of networks
        $N>2$, comparing each network against the mean utilization of the
        other $N-1$ networks. We prove that the formula reduces exactly to
        the base pairwise signal when $N=2$, and confirm the reduction
        numerically at absolute tolerance $10^{-7}$.
  \item \textbf{Distance-aware balancing.} We ground each network at a real
        geographic coordinate and weight each neighbor's contribution to the
        balancing signal by a Haversine-distance exponential decay, placing
        the collaborative pricing loop on the same footing as classical
        gravity models of spatial interaction~\cite{huff1964}.
  \item \textbf{Renewable net-load integration.} We derive physics-based
        solar-photovoltaic and wind generation estimates directly from the
        real meteorological columns already present in the public dataset,
        and reformulate the pricing and demand-response equations in terms
        of renewable-adjusted net load with a $5\%$ baseload floor.
  \item \textbf{Validation and inheritance analysis.} On a four-network
        testbed grounded in real Tetouan, Morocco telemetry, we show that
        the combined framework (PVB\_FULL) reduces the maximum--minimum
        cross-network utilization imbalance by $19.03\%$ relative to
        uncoordinated pricing and by $15.28\%$ relative to the original
        pairwise formulation, with the ranking preserved across four
        seasonal windows, five inference-noise seeds, and four DDPG
        training seeds. A four-way ablation isolates the marginal
        contribution of each extension, and a direct re-analysis of the
        base paper's own two-network data reveals that a structural
        EV peak-uniformity property is inherited from its independent
        min-max normalization.
\end{enumerate}

The remainder of the paper is organized as follows.
Section~\ref{sec:rw} reviews related work. Section~\ref{sec:prelim}
formalizes the problem and the base formulations. Section~\ref{sec:framework}
presents the three extensions and the combined architecture.
Section~\ref{sec:setup} describes the testbed and evaluation protocol.
Section~\ref{sec:results} reports the results, including the
inherited-limitation analysis. Section~\ref{sec:disc} discusses the distance
decay, the inherited uniformity limitation, and threats to validity.
Section~\ref{sec:conc} concludes and outlines future work.

% ======================================================================
\section{Related Work}\label{sec:rw}

\subsection{Time-of-Use and Dynamic Pricing}
TOU tariffs remain the most widely deployed EV price signal, yet their
empirical record is mixed. In a large-scale residential field experiment,
Bailey \emph{et al.} observed that TOU pricing shifted charging out of peak
periods but created larger simultaneous shadow peaks, increasing local
capacity-exceedance risk~\cite{bailey2025}. Complementing this with real
distribution-transformer telemetry, Samende \emph{et al.} demonstrated a
``herding effect'' in which uncoordinated smart chargers responding to one
discounted window produce a new peak that, above roughly $20$--$30\%$ EV
penetration, exceeds the traditional evening
peak~\cite{samende2025}. Simulation studies corroborate the effect, with
reported induced peaks of up to $20\%$ when most charging begins immediately
after on-peak periods end~\cite{jones2022}. On the economics side, Faruqui
and Sergici's survey of fifteen pilots quantifies the peak-demand response
achievable under different tariff granularities~\cite{faruqui2010}, and
Borenstein's long-run analysis quantifies the efficiency ceiling of coarse
two-part tariffs relative to real-time retail prices~\cite{borenstein2005}.
These works imply that tariff design should adapt continuously and be
spatially differentiated--the direction pursued here.

\subsection{EV Charging Control and Distribution-Grid Impact}
Charging management has been treated as a network problem for over a
decade. Clement-Nyns \emph{et al.} showed that uncoordinated PHEV charging
raises residential losses and voltage deviations, and that coordinated
charging mitigates both~\cite{clement2010}. Richardson \emph{et al.}
compared local and centralized charging strategies on low-voltage
feeders~\cite{richardson2012tsg}, and showed earlier that LP-based optimal
charging maximizes the EV capacity a feeder can host without
reinforcement~\cite{richardson2012tpwrs}. O'Connell \emph{et al.}
introduced a rolling multi-period optimization that re-solves charging
plans hourly in response to forecast
updates~\cite{oconnell2014}. Analytically, Gan, Topcu, and Low proved that a
decentralized price-update protocol converges to an aggregate-valley-filling
charging profile~\cite{gan2013}, Ma \emph{et al.} established a Nash
equilibrium for decentralized charging of large vehicle
populations~\cite{ma2013}, and Callaway and Hiskens characterized the
controllability of aggregate electric loads including
EVs~\cite{callaway2011}. At the distribution-operation level, the
load-balancing view of feeder operation goes back to the reconfiguration
framework of Baran and Wu~\cite{baran1989}. All of these works treat
imbalance as a cross-feeder, cross-time problem at its core--the same view
adopted by the collaborative pricing framework we generalize.

\subsection{Deep Reinforcement Learning for Pricing and Demand Response}
DDPG~\cite{lillicrap2016} extends the deep Q-learning paradigm
\cite{mnih2015} to continuous action spaces and, together with its
off-policy entropy-maximizing and on-policy successors
\cite{haarnoja2018,schulman2017}, defines the practical toolbox for
continuous price and dispatch control. Survey evidence shows that RL-based
demand response is effective but remains dominated by single-agent,
single-network settings with demand-independent prices--precisely the gaps
that multi-agent price coordination is meant to fill~\cite{vazquez2019}.
Applied to EV charging, DDPG has been reported for aggregator tariff
optimization and grid-load smoothing~\cite{qiu2020tia,liu2021access};
station-level revenue and response-aware pricing \cite{zhao2022tte},
\cite{lai2022tsg}; public-station pricing in coupled transportation-power
systems~\cite{cui2021tsg}; charging control with dynamic user behavior
\cite{yan2021tsg}; multi-agent charge/discharge scheduling
\cite{aljafari2023}; battery-assisted fast-charging hub management
\cite{paudel2023}; coordinated energy trading and flexibility provision in
local markets~\cite{ye2022}; and competitive fast-charging-station pricing
games~\cite{yang2024}. None of these coordinates the tariffs of one network
with the grid state of multiple \emph{other distribution networks} while
retaining a closed-form demand-response baseline; the pairwise
collaborative formulation~\cite{lepolesa2026} remains the closest
predecessor and is the object extended here.

\subsection{Spatial Choice, Distance, and Charging Behavior}
EV owners choose charging locations as a spatial-choice problem. The
gravity paradigm of Huff models the probability of patronizing a site as an
increasing function of attractiveness and a decreasing function of
distance, formalizing the ``friction of distance'' in
consumer behavior~\cite{huff1964}. Charging-network planning applies these
ideas directly: maximal-covering and accessibility models place stations
where they serve the largest nearby demand at acceptable travel
cost~\cite{frade2011}. At the control level, DRL formulations that embed
user mobility behavior show that ignoring the coupling between tariff and
travel choice degrades both demand shaping and
profitability~\cite{yan2021tsg}. These results motivate the distance-decay
weighting in Section~\ref{sec:framework}: if willingness to travel for a
cheaper tariff decays with distance, then a distant network's influence on a
balancing signal should decay accordingly.

\subsection{Load Forecasting}
All balancing strategies require forecasts of conventional load before
pricing can act. Gradient-boosted decision trees--XGBoost~\cite{chen2016}
and LightGBM~\cite{ke2017}--and random forests~\cite{breiman2001} dominate
tabular load-forecasting benchmarks, while recurrent architectures such as
LSTMs~\cite{hochreiter1997} capture sequential structure; probabilistic
treatments and systematic evaluation practices are emphasized in recent
tutorials~\cite{hong2016,wang2016recency}. The base framework uses boosting
for this purpose~\cite{lepolesa2026}, and we retain that choice, using only
the forecasted load as input to the demand-response and pricing stages.

% ======================================================================
\section{Problem Formulation and Preliminaries}\label{sec:prelim}

\subsection{Network Model and Notation}
Let $\mathcal{N}=\{1,\dots,N\}$ denote the set of distribution networks,
each with rated capacity $C_{\max,i}$ and time-varying load $P_i(t)$. The
instantaneous utilization of network $i$ is
\begin{equation}\label{eq:util}
  \util_i(t)=\frac{P_i(t)}{C_{\max,i}},\qquad \util_i(t)\in(0,1].
\end{equation}
Primary performance indicators used throughout are the time-averaged
maximum--minimum utilization spread
\begin{equation}\label{eq:imb}
  \bar{I}=\frac{1}{T}\sum_{t=1}^{T}\Bigl(
    \max_{i\in\mathcal{N}}\util_i(t)-\min_{i\in\mathcal{N}}\util_i(t)\Bigr),
\end{equation}
and the time-averaged cross-network utilization standard deviation
\begin{equation}\label{eq:sid}
  \bar{S}=\frac{1}{T}\sum_{t=1}^{T}s(t),\qquad
  s(t)=\sqrt{\frac{1}{N-1}\sum_{i=1}^{N}\bigl(
    \util_i(t)-\overline{\util}(t)\bigr)^2},
\end{equation}
where $\overline{\util}(t){=}\tfrac{1}{N}\sum_i\util_i(t)$. Table~\ref{tab:not}
summarizes the symbols used in the remainder of the paper.

\begin{table}[!t]
\caption{Principal Notation.}
\label{tab:not}
\centering
\begin{tabular}{ll}\toprule
Symbol & Meaning \\ \midrule
$\mathcal{N}$, $N$ & Network set and cardinality \\
$C_{\max,i}$ & Rated capacity of network $i$ \\
$P_i(t)$, $P_{\mathrm{net},i}(t)$ & Conventional and net load of network $i$ \\
$\util_i(t)$ & Utilization of network $i$ \\
$\D U_i(t)$ & Balancing signal of network $i$ \\
$d_{ij}$, $W_{ij}$ & Haversine distance and willingness decay \\
$\mu_i$ & Per-network EV demand scale \\[2pt]
$\varepsilon(t)$ & Gaussian EV-demand perturbation \\ \bottomrule
\end{tabular}
\end{table}

\subsection{Load Forecasting Module}
Conventional load $P^{\mathrm{conv}}_i(t)$ is predicted hour ahead from
lagged load and weather covariates using gradient-boosted trees trained as
in the base framework~\cite{lepolesa2026,chen2016}. Forecast quality is
assessed with the standard point-forecast metrics
\begin{align}
  \mathrm{MAE}_i&=\frac{1}{T}\sum_{t=1}^{T}\bigl|
    P_i^{\mathrm{conv}}(t)-\hat P_i(t)
    \bigr|,\label{eq:mae}\\
  \mathrm{RMSE}_i&=\sqrt{\frac{1}{T}\sum_{t=1}^{T}\bigl(
    P_i^{\mathrm{conv}}(t)-\hat P_i(t)\bigr)^2},\label{eq:rmse}
\end{align}
where $\hat P_i(t)$ denotes the forecast. The forecast column is the common
input shared by every strategy compared in Section~\ref{sec:results}, so
forecast error affects all baselines equally.

\subsection{Closed-Form EV Demand Response (Base Formulation)}
The base paper models EV charging demand as a direct, closed-form response
to the network state, with no trained demand model
\cite{lepolesa2026}. Let $\mu_i$ be the per-network EV demand scale
($\mu_i=0.5\max_t\{D_i^{\mathrm{unc}}(t)\}$, half the unbalanced EV-demand
peak, matching the reference implementation), and let
$\norm{x}=\tfrac{x-\min x}{\max x-\min x}$ denote independent per-network
min-max normalization applied to the entire $T$-sample. Under the
peak-valley (PV) strategy the EV demand is
\begin{equation}\label{eq:drpv}
  D_i^{\mathrm{PV}}(t)=\mu_i\bigl(1-\norm{P_i(t)}\bigr)+\varepsilon(t),
\end{equation}
and under the peak-valley-with-balancing (PVB) strategy
\begin{equation}\label{eq:drpvb}
  D_i^{\mathrm{PVB}}(t)=\frac{\mu_i}{2}\Bigl[
    \bigl(1-\norm{P_i(t)}\bigr)
    +\bigl(1-\norm{\D U_i(t)}\bigr)\Bigr]+\varepsilon(t),
\end{equation}
where $\D U_i(t)$ is the balancing signal defined below and
$\varepsilon(t)$ is a small Gaussian perturbation representing
prediction/behavior noise. The resulting grid load is
$P_i(t)+D_i(t)$, and the per-network PAR over the horizon is
\begin{equation}\label{eq:par}
  \mathrm{PAR}_i=
  \frac{\max_{t\in[1,T]}\bigl[P_i(t)+D_i(t)\bigr]}
  {\frac{1}{T}\sum_{t=1}^{T}\bigl[P_i(t)+D_i(t)\bigr]}.
\end{equation}

\subsection{Base Two-Network Balancing}
For the two-network case $N{=}2$, the balancing signal of the base
formulation is the clipped pairwise utilization difference
\begin{equation}\label{eq:pairwise}
  \D U_{i}(t)=\max\bigl(\util_i(t)-\util_j(t),\,0\bigr),\qquad j\ne i.
\end{equation}
Network $i$ is induced to export demand to $j$ whenever it is more loaded
than $j$; the clip prevents negative balancing signals from reversing the
incentive. This pairwise construction is the object generalized in
Section~\ref{sec:framework}.

\subsection{Problem Statement}
We seek, for arbitrary $N\ge2$, pricing strategies $p_i(t)$ that minimize
the cross-network imbalance $\bar I$ (\ref{eq:imb}) while preserving per-network
valley-filling. Three subproblems are posed, corresponding to the three
gaps in the base formulation.
\begin{itemize}
  \item \textbf{P1 (Scalability).} Construct a balancing signal
        $\D U_i(t)$ for arbitrary $N$ that reduces exactly to
        (\ref{eq:pairwise}) when $N{=}2$.
  \item \textbf{P2 (Spatial friction).} Incorporate real geographic
        distance so that the influence of network $j$ on the balancing
        signal of network $i$ decays with $d_{ij}$.
  \item \textbf{P3 (Renewables).} Reformulate $P_i(t)$ in the demand
        equations using a renewable-adjusted net load that retains the
        closed-form structure of (\ref{eq:drpv})--(\ref{eq:drpvb}).
\end{itemize}

% ======================================================================
\section{Proposed Framework}\label{sec:framework}

\subsection{System Architecture}
Fig.~\ref{fig:arch} shows the overall architecture. For each network a
forecast module produces $\hat P_i^{\mathrm{conv}}(t)$; a renewable module
displaces it into net load $P_{\mathrm{net},i}(t)$
(Section~\ref{sec:renewables}); the generalized, distance-weighted
balancing signal $\D U_{i,\mathrm{dist}}(t)$ (Sections~\ref{sec:gen} and
\ref{sec:decay}) is constructed from the net-load utilizations; and the
closed-form EV demand---and, for dynamic pricing, the collaborative DDPG
tariff $p_i(t)$---closes the loop.

\begin{figure}[!t]
\centering
\begin{tikzpicture}[
  >=Stealth,
  font=\scriptsize,
  inp/.style={draw=black!65, rounded corners=2pt, align=center, inner sep=4pt,
              fill=black!5, minimum width=2.6cm},
  proc/.style={draw=black!65, rounded corners=2pt, align=center, inner sep=5pt,
               fill=black!13, minimum width=5.9cm},
  out/.style={draw=black!65, rounded corners=2pt, align=center, inner sep=5pt,
              fill=black!4, minimum width=5.9cm},
  arr/.style={->, >=Stealth[length=2.3mm,width=2.0mm], line width=0.7pt},
  stg/.style={font=\tiny\bfseries, gray!50}
]
\node[font=\tiny\itshape, gray!70] at (0,3.4)
  {For each network $i\in\{1,\dots,N\}$};
\node[inp] (fc) at (-1.9,2.35) {Conventional load\\forecast (GBM)};
\node[inp] (ren) at (1.9,2.35) {Renewable generation\\(PV 18\%, wind law)};
\node[proc] (net) at (0,0.85) {Net load and utilization\\$u_i(t)=P_{\mathrm{net},i}(t)/C_{\max,i}$};
\node[proc] (bal) at (0,-0.95) {Distance-aware balancing signal\\$\D U_{i,\mathrm{dist}}(t)$\ (arbitrary $N$, Haversine $W_{ij}$)};
\node[proc] (dem) at (0,-2.75) {Closed-form EV demand and tariff\\$D_i^{\mathrm{FULL}}(t)$\ under
  $p_i{=}p_{\mathrm{conv}}{+}\tfrac12(a_{\mathrm{P1},i}{+}a_{\mathrm{P2},i})$};
\node[out] (met) at (0,-4.55) {Reporting\\$\bar I$, $\bar S$, per-network PAR, paired $t$-test};
\draw[arr] (fc) -- (net);
\draw[arr] (ren) -- (net);
\draw[arr] (net) -- (bal);
\draw[arr] (bal) -- (dem);
\draw[arr] (dem) -- (met);
\node[stg] at (-4.05,0.9) {1};
\node[stg] at (-4.05,-0.9) {2};
\node[stg] at (-4.05,-2.7) {3};
\node[stg] at (-4.05,-4.5) {4};
\end{tikzpicture}
\caption{System architecture of the generalized framework. Stage~1:
forecasting and on-site renewables reduce the feeder load to a renewable
net load $P_{\mathrm{net},i}(t)$. Stage~2: a distance-weighted balancing
signal aggregates all $N$ networks. Stage~3: closed-form demand response is
priced by collaborative DDPG tariffs. Stage~4: metrics are reported. Only
the balancing signal and the load variable change between PV, PVB\_OLD,
PVB\_NEW, and PVB\_FULL.}
\label{fig:arch}
\end{figure}

Forecast, balancing, and demand modules are computed from the same
closed-form equations for every compared strategy; only the construction of
the balancing signal and the load variable changes between PV, PVB\_OLD,
PVB\_NEW, and PVB\_FULL. This guarantees that all differences reported in
Section~\ref{sec:results} are attributable to the three extensions alone.

\subsection{Contribution 1: Generalized $N$-Network Balancing}\label{sec:gen}
For arbitrary $N\ge2$, network $i$ is compared against the mean utilization
of all other networks rather than a single partner:
\begin{equation}\label{eq:gen}
  \D U_i(t)=\max\Bigl(\util_i(t)-\frac{1}{N-1}\sum_{j\ne i}\util_j(t),\,
    0\Bigr).
\end{equation}

\begin{proposition}[Backward compatibility]\label{prop:n2}
For $N=2$, (\ref{eq:gen}) reduces exactly to the base pairwise signal
(\ref{eq:pairwise}).
\end{proposition}
\begin{IEEEproof}
With $N=2$, the single-neighbor mean in (\ref{eq:gen}) collapses to
$\tfrac{1}{2-1}\util_j(t)=\util_j(t)$ for the one network $j\ne i$, so
$\D U_i(t)=\max(\util_i(t)-\util_j(t),0)$, identical to
(\ref{eq:pairwise}).
\end{IEEEproof}

The reduction is additionally verified in code: the companion notebook
asserts, with $\texttt{np.allclose}$, that the generalized function and the
base pairwise function produce identical balancing signals at absolute
tolerance $10^{-7}$ on the canonical inputs (Appendix~\ref{app:equiv}).
This dual verification guards against the algebraic argument alone.

\subsection{Contribution 2: Spatial Distance Decay}\label{sec:decay}
The base formulation implicitly assumes infinite willingness to travel to
any other network. We replace that assumption with a real-geography,
gravity-style decay~\cite{huff1964}. Each network is placed at the centroid
$(\phi_i,\lambda_i)$ of its zone, and the great-circle distance between two
networks is computed with the Haversine formula
\begin{equation}\label{eq:hav}
  d_{ij}=2R\arcsin\!\Bigl(\sqrt{
    \sin^2\!\tfrac{\D\phi_{ij}}{2}
    +\cos\phi_i\cos\phi_j\,\sin^2\!\tfrac{\D\lambda_{ij}}{2}}\Bigr),
\end{equation}
with $R=6371~\mathrm{km}$. Driver willingness to shift charging from $i$ to
$j$ decays exponentially with distance,
\begin{equation}\label{eq:w}
  W_{ij}=\exp(-\alpha\,d_{ij}),\qquad W_{ii}=0,\qquad \alpha=0.15\,\mathrm{km^{-1}},
\end{equation}
and the balancing signal replaces the unweighted mean in
(\ref{eq:gen}) with the row-normalized weighted mean
\begin{equation}\label{eq:dist}
  \D U_{i,\mathrm{dist}}(t)=\max\Bigl(\util_i(t)-\sum_{j\ne i}
   \frac{W_{ij}}{\sum_{k\ne i}W_{ik}}\,\util_j(t),\,
   0\Bigr).
\end{equation}
For $\alpha\to0$, (\ref{eq:dist}) degenerates smoothly to the unweighted
mean of (\ref{eq:gen}); for $\alpha\to\infty$, it degenerates to the
nearest-neighbor pairwise signal. Table~\ref{tab:dist} lists the pairwise
distances used here; all are typical of an intra-city topology.

\begin{table}[!t]
\footnotesize
\setlength{\tabcolsep}{5pt}
\caption{Haversine Distance Matrix Between Zone Centroids (km).}
\label{tab:dist}
\centering
\begin{tabular}{lcccc}\toprule
 & Res. & Com. & Ind. & Inst. \\ \midrule
Residential    & 0    & 1.16 & 5.52 & 1.54 \\
Commercial     & 1.16 & 0    & 5.15 & 2.28 \\
Industrial     & 5.52 & 5.15 & 0    & 7.06 \\
Institutional  & 1.54 & 2.28 & 7.06 & 0    \\ \bottomrule
\end{tabular}
\end{table}

\subsection{Contribution 3: Renewable Net-Load Integration}\label{sec:renewables}
The public dataset supplies hourly global and diffuse irradiance components
and wind speed alongside conventional load
\cite{salam2018}. Rather than importing an external scenario, we derive
renewable generation from these measured columns. Solar photovoltaic
generation is a linear irradiance-PV-aperture model,
\begin{equation}\label{eq:solar}
  P_{\mathrm{solar},i}(t)=G(t)\,A_{\mathrm{PV},i}\,\eta_{\mathrm{PV}},
\end{equation}
where $G(t)$ is the total (direct$+$diffuse) in-plane irradiance
($\mathrm{W/m^2}$), $A_{\mathrm{PV},i}$ the rooftop aperture assumed for
network $i$, and $\eta_{\mathrm{PV}}=18\%$ the panel efficiency. Wind
generation follows a clipped cubic power curve
\begin{equation}\label{eq:wind}
P_{\mathrm{wind},i}(t)=
\begin{cases}
0, & v(t)<v_{\mathrm{cutin}}\ \text{or}\ v(t)\ge v_{\mathrm{cutout}},\\[2pt]
\tfrac12\rho A_{\mathrm{rotor}}C_p\,v(t)^3\,n_i,\;
    & v_{\mathrm{cutin}}\le v(t)<v_{\mathrm{rated}},\\[2pt]
P_{\mathrm{rated}}\,n_i,\; & v_{\mathrm{rated}}\le v(t)<v_{\mathrm{cutout}},
\end{cases}
\end{equation}
with air density $\rho=1.225~\mathrm{kg/m^3}$, power coefficient
$C_p=0.35$, $n_i$ turbines per network, cut-in
$v_{\mathrm{cutin}}=1.5\,\mathrm{m/s}$, rated
$v_{\mathrm{rated}}=9.0\,\mathrm{m/s}$, and cut-out
$v_{\mathrm{cutout}}=20.0\,\mathrm{m/s}$.
The renewable-adjusted net load, subject to a $5\%$ baseload floor that
prevents an unrealistic zero-load condition during peak midday generation,
is
\begin{equation}\label{eq:net}
  P_{\mathrm{net},i}(t)=\max\Bigl(0.05\,P_{\mathrm{conv},i}(t),\;
    P_{\mathrm{conv},i}(t)-P_{\mathrm{solar},i}(t)-P_{\mathrm{wind},i}(t)\Bigr).
\end{equation}
High PV penetration is known to reshape feeder net-load and voltage
envelopes in precisely this way~\cite{hoke2013}, motivating the
net-load-centric reformulation. Fig.~\ref{fig:netload} illustrates the
conventional versus net-load profiles for the four test networks.

\paperfig[\columnwidth]{fig3_net_load_profiles.png}%
{Conventional versus renewable-adjusted net load for the four test networks
on the canonical 24-h window.}%
{fig:netload}

\subsection{EV Demand Response under the Combined Framework}
PVB\_FULL is obtained by substituting net load for conventional load and the
distance-weighted net-load signal for $\D U_i$ in the closed-form PVB
equation:
\begin{equation}\label{eq:full}
\begin{aligned}
D_i^{\mathrm{FULL}}(t)=\frac{\mu_i}{2}\Bigl[&
  \bigl(1-\norm{P_{\mathrm{net},i}(t)}\bigr)\\
  &+{\bigl(1-\norm{\D U_{i,\mathrm{dist}}(t)}\bigr)}\Bigr]
  +\varepsilon(t),
\end{aligned}
\end{equation}
with $\D U_{i,\mathrm{dist}}$ computed from net-load utilizations.
Equations (\ref{eq:drpv}) and (\ref{eq:drpvb}) remain the definitions of the
PV and PVB\_OLD baselines, so the four strategies differ only through the
extensions.

\subsection{Collaborative DDPG Dynamic Pricing}\label{sec:framework-ddpg}
For dynamic retail pricing each network is assigned two DDPG
actor--critic pairs~\cite{lillicrap2016}, following the base
architecture~\cite{lepolesa2026}. Agent P1 targets the network's own
normalized net load with a mean-absolute-error (MAE) reward; agent P2
targets the normalized distance-weighted balancing signal with an MAE
reward. Actors use a $24{-}64{-}64{-}24$ feedforward network with ReLU
activations and a scaled-tanh output bounded to $[0,1]$; critics use
$48{-}64{-}64{-}1$ networks over the concatenated state--action vector, and
all networks are optimized with Adam~\cite{kingma2015}. The final hourly
tariff is the convex combination
\begin{equation}\label{eq:price}
  p_i(t)=p_{\mathrm{conv}}(t)
       +\tfrac12\,a_{\mathrm{P1},i}(t)
       +\tfrac12\,a_{\mathrm{P2},i}(t),
\end{equation}
where $p_{\mathrm{conv}}(t)$ is the conventional retail tariff.

\subsubsection*{Evaluation Protocol and the Role of DDPG}
Consistent with the base paper's reported methodology (verified against its
public reference code), the headline utilization metrics in
Section~\ref{sec:results} are computed from the closed-form demand equations
above; the DDPG actions enter only through the tariff $p_i(t)$ used to
construct the demand. The learned pricing policy is evaluated separately: it
is retrained from scratch over four training seeds and its run-to-run
tariff variability is reported (Section~\ref{sec:ddpg-rob}). This two-tier
separation is a faithful match to the base method~\cite{lepolesa2026}, not an
implementation shortcut.

% ======================================================================
\section{Experimental Setup}\label{sec:setup}

\subsection{Testbed and Data}
Residential and commercial load profiles are the real hourly series used by
the base paper: $8{,}736$ hours of Tetouan, Morocco, distribution telemetry
(residential feeder) together with a real commercial feeder series and the
base paper's EV-demand columns, all from the public Tetouan smart-grid
dataset~\cite{salam2018}. Industrial and institutional profiles are
synthesized by re-weighting the real residential series with distinct
diurnal envelopes--a continuous day-shift plateau from $06{:}00$ to $18{:}00$
for industrial and a bimodal university/hospital schedule
($08{:}00$--$12{:}00$, $14{:}00$--$17{:}00$) for institutional--following
the same class of profile construction the base paper applies to its
synthetic EV demand. Rated capacities are $1.167$, $1.0$, $1.5$, and $0.8$
MW for the residential, commercial, industrial, and institutional networks,
respectively. Zone centroids are placed within the city of Tetouan, giving
the pairwise distances of Table~\ref{tab:dist}. Fig.~\ref{fig:pipeline}
shows the end-to-end evaluation pipeline.

\begin{figure}[!t]
\centering
\begin{tikzpicture}[
  >=Stealth,
  font=\scriptsize,
  st/.style={draw=black!65, rounded corners=2pt, align=center, inner sep=5pt,
             minimum width=6.6cm},
  arr/.style={->, >=Stealth[length=2.3mm,width=2.0mm], line width=0.7pt,
              shorten >=2pt}
]
\node[st,fill=black!5] (s0) at (0,2.9)
  {Stage 0 --- data acquisition\\Tetouan 10-min telemetry, EV demand,
  solar/wind columns};
\node[st,fill=black!13] (s1) at (0,1.6)
  {Stage 1 --- conventional load forecast\\GBM/XGBoost day-ahead
  $\hat P_i^{\mathrm{conv}}(t)$};
\node[st,fill=black!13] (s2) at (0,0.3)
  {Stage 2 --- renewable net load\\$P_{\mathrm{net},i}(t)$ with $5\%$
  baseload floor};
\node[st,fill=black!13] (s3) at (0,-1.0)
  {Stage 3 --- generalized balancing signal\\$\D U_{i,\mathrm{dist}}(t)$,
  $\alpha=0.15~\mathrm{km^{-1}}$};
\node[st,fill=black!4] (s4) at (0,-2.3)
  {Stage 4 --- evaluation\\closed-form $D_i(t)$ $+$ DDPG
  $p_i(t)$ $\to$ $\bar I$, $\bar S$, PAR, paired $t$-test};
\draw[arr] (s0)--(s1);
\draw[arr] (s1)--(s2);
\draw[arr] (s2)--(s3);
\draw[arr] (s3)--(s4);
\end{tikzpicture}
\caption{Evaluation and training pipeline from raw telemetry to the
reported metrics (Stage 0--4 of the companion notebook).}
\label{fig:pipeline}
\end{figure}

\subsection{Evaluation Protocol}
We report $\bar I$ (\ref{eq:imb}), $\bar S$ (\ref{eq:sid}), and per-network PAR
(\ref{eq:par}). Strategies are evaluated on the canonical $24$-hour window
at hour $576$ of the annual series (the base paper's own test index) and,
for the seasonal study, on three additional windows from spring
(hr~$2400$), summer (hr~$4800$), and autumn (hr~$7200$), under two EV-fleet
scaling assumptions (fixed annual $\mu_i$ vs.\ per-window recomputed
$\mu_i$). Five fixed inference-noise seeds
$\{42,101,2024,777,999\}$ quantify closed-form sensitivity to
$\varepsilon(t)$; these are separate from the four DDPG training seeds
$\{42,7,123,2025\}$, since the two seed sets control distinct sources of
variability. All numbers reported in the tables are the canonical outputs of
the companion notebook (files in \texttt{data/}).

% ======================================================================
\section{Results and Discussion}\label{sec:results}

\subsection{Headline Strategy Benchmark}
Table~\ref{tab:bench} compares PV, PVB\_OLD (the pairwise formula applied to
four networks), PVB\_NEW (generalized $N$-network balancing alone,
Contribution~1), and PVB\_FULL (all three contributions). PVB\_FULL reduces
the maximum--minimum imbalance by $19.03\%$ relative to PV and by $15.28\%$
relative to PVB\_OLD; the time-averaged cross-network utilization standard
deviation falls from $0.1152$ to $0.0901$, a $21.81\%$ reduction. The
observed standard deviations over the five inference seeds are two orders of
magnitude below the strategy differences, so the ranking is robust to the
closed-form demand noise. The per-hour imbalance envelope is plotted in
Fig.~\ref{fig:imbalance}.

\begin{table}[!t]
\footnotesize
\setlength{\tabcolsep}{4pt}
\caption{Headline Strategy Benchmark (mean${\pm}$std over 5 inference seeds,
canonical 24-h window).}
\label{tab:bench}
\centering
\begin{tabular}{lccc}\toprule
Strategy & $\bar I$ & $\bar S$ & Mean PAR \\ \midrule
PV        & $0.2729{\pm}0.0008$ & $0.1152{\pm}0.0003$ & $1.344{\pm}0.010$ \\
PVB\_OLD  & $0.2608{\pm}0.0001$ & $0.1116{\pm}0.0000$ & $1.335{\pm}0.001$ \\
PVB\_NEW  & $0.2544{\pm}0.0001$ & $0.1044{\pm}0.0000$ & $1.378{\pm}0.001$ \\
PVB\_FULL & $0.2210{\pm}0.0001$ & $0.0901{\pm}0.0000$ & $1.416{\pm}0.002$ \\ \bottomrule
\end{tabular}
\end{table}

\paperfig[\columnwidth]{fig4_imbalance_by_strategy.png}%
{Time series of the per-hour maximum--minimum utilization envelope for the
four strategies, averaged over the five inference seeds with the seed
min--max range shaded. Every point reproduces the Stage-4 canonical
evaluation.}%
{fig:imbalance}

\subsection{Ablation Study}
Table~\ref{tab:abl} isolates each extension's marginal contribution by
evaluating four nested configurations. Renewable net-load integration
(Contribution~3) accounts for nearly all of the combined gain
($-12.88\%$), while spatial distance decay (Contribution~2) contributes a
modest but consistent $-0.30\%$, as visualized in Fig.~\ref{fig:ablation}.
The asymmetry is expected at intra-city
distances of $1.16$--$7.06$~km: transit time is short relative to the
several-hundred-kilowatt midday load removed by on-site generation.

\begin{table}[!t]
\footnotesize
\setlength{\tabcolsep}{4pt}
\caption{Four-Way Ablation Study ($\bar I$, 5-seed mean).}
\label{tab:abl}
\centering
\begin{tabular}{lccc}\toprule
Configuration & $\bar I$ & Marginal $\D$ & Marginal \% \\ \midrule
(a) Gap 1 only            & 0.2544 & ---    & --- \\
(b) Gap 1 $+$ Gap 2       & 0.2536 & $-0.0008$ & $-0.30\%$ \\
(c) Gap 1 $+$ Gap 3       & 0.2216 & $-0.0328$ & $-12.88\%$ \\
(d) All three (PVB\_FULL) & 0.2210 & $-0.0334$ & $-13.13\%$ \\ \bottomrule
\end{tabular}
\end{table}

\paperfig[\columnwidth]{fig5_ablation_contributions.png}%
{Ablation: $\bar I$ for configurations (a)--(d), with the marginal change
relative to (a) shown above each bar. Values are the canonical 5-seed
outputs of Table~\ref{tab:abl}.}%
{fig:ablation}

\subsection{Statistical Significance}
A paired $t$-test comparing PVB\_OLD and PVB\_FULL over the five
inference-noise seeds uses the statistic
\begin{equation}\label{eq:ttest}
  t=\frac{\bar d}{s_d/\sqrt{n}},\qquad
  d_k=\bar I_k^{\mathrm{PVB\_OLD}}-\bar I_k^{\mathrm{PVB\_FULL}},
  \quad k=1,\dots,5,
\end{equation}
and yields a mean reduction of $\bar d=0.039840$
(95\% CI $[0.039832,0.039849]$, $t=12653.6$, $p=2.34\times10^{-16}$). We
emphasize what this interval means: it characterizes sensitivity to the
small Gaussian perturbation $\varepsilon(t)$ applied to an otherwise
deterministic closed-form formula, i.e., consistency across demand-noise
realizations. It does \emph{not} represent inter-grid stochastic
variability, which the closed-form model does not simulate; DDPG training
variability is treated separately in Section~\ref{sec:ddpg-rob}.

\subsection{Seasonal Generalization}
Tables~\ref{tab:season} and~\ref{tab:seasonb} report the two EV-fleet
scaling assumptions (Case~A: fixed annual $\mu_i$; Case~B: $\mu_i$
recomputed from each season's conventional peak). The ordering
$\mathrm{PV}>\mathrm{PVB\_OLD}>\mathrm{PVB\_NEW}>\mathrm{PVB\_FULL}$ holds
in all eight evaluations (Fig.~\ref{fig:seasonal}), with reductions of
$16.5$--$19.0\%$ versus PV and
$11.4$--$15.3\%$ versus PVB\_OLD across the full set.

\begin{table}[!t]
\footnotesize
\setlength{\tabcolsep}{4pt}
\caption{Seasonal Generalization of $\bar I$ --- Case A (fixed annual
$\mu_i$).}
\label{tab:season}
\centering
\begin{tabular}{lcccc}\toprule
Season & PV & PVB\_OLD & PVB\_NEW & PVB\_FULL \\ \midrule
Winter & 0.2729 & 0.2608 & 0.2544 & 0.2210 \\
Spring & 0.2643 & 0.2491 & 0.2417 & 0.2206 \\
Summer & 0.3209 & 0.3016 & 0.2945 & 0.2645 \\
Autumn & 0.2556 & 0.2418 & 0.2324 & 0.2088 \\ \bottomrule
\end{tabular}
\end{table}

\begin{table}[!t]
\footnotesize
\setlength{\tabcolsep}{4pt}
\caption{Seasonal Generalization of $\bar I$ --- Case B (per-window
recomputed $\mu_i$).}
\label{tab:seasonb}
\centering
\begin{tabular}{lcccc}\toprule
Season & PV & PVB\_OLD & PVB\_NEW & PVB\_FULL \\ \midrule
Winter & 0.2729 & 0.2608 & 0.2544 & 0.2210 \\
Spring & 0.2672 & 0.2522 & 0.2438 & 0.2230 \\
Summer & 0.3327 & 0.3114 & 0.3041 & 0.2754 \\
Autumn & 0.2588 & 0.2456 & 0.2357 & 0.2126 \\ \bottomrule
\end{tabular}
\end{table}

\paperfig[\columnwidth]{fig6_seasonal_results.png}%
{Seasonal generalization of $\bar I$ for Case A (fixed annual $\mu_i$) and
Case B (per-season $\mu_i$); each group shows PV, PVB\_OLD, PVB\_NEW, and
PVB\_FULL, matching Tables~\ref{tab:season} and~\ref{tab:seasonb}.}%
{fig:seasonal}

\subsection{DDPG Pricing-Policy Robustness}\label{sec:ddpg-rob}
Because the headline metrics are closed-form, DDPG training variability is
evaluated directly on the learned tariff. All eight collaborative agents
were retrained from scratch over the four training seeds
$\{42,7,123,2025\}$. Table~\ref{tab:ddpg} reports that learned tariffs show
$6.7$--$13.0\%$ relative standard deviation across training runs while
consistently taxing morning/evening peak hours and discounting midday hours
that coincide with peak solar availability in every run
(Fig.~\ref{fig:ddpg}).

\begin{table}[!t]
\footnotesize
\setlength{\tabcolsep}{5pt}
\caption{DDPG Tariff Variability Across Four Training Seeds.}
\label{tab:ddpg}
\centering
\begin{tabular}{lcccc}\toprule
Network & Mean tariff & Mean seed std & Max seed std & Rel.\ var. \\ \midrule
Residential    & 0.950 & 0.123 & 0.275 & $13.0\%$ \\
Commercial     & 0.938 & 0.113 & 0.320 & $12.1\%$ \\
Industrial     & 1.030 & 0.094 & 0.247 & $9.1\%$ \\
Institutional  & 0.953 & 0.064 & 0.191 & $6.7\%$ \\ \bottomrule
\end{tabular}
\end{table}

\paperfig[\columnwidth]{fig7_ddpg_prices.png}%
{Representative learned dynamic tariffs $p_i(t)$ of the canonical policy
for the four networks over the 24-h horizon, together with the conventional
tariff.}%
{fig:ddpg}

\subsection{Peak-to-Average Ratio Decomposition}\label{sec:par}
PAR increases slightly under PVB\_FULL despite falling peaks
(Fig.~\ref{fig:par}). Table~\ref{tab:par} decomposes the cause per network:
absolute peak demand falls
in every network, but midday solar generation depresses the daily mean load
(the PAR denominator) by $5.7$--$10.8\%$, more than it reduces the
non-solar residual peak (the numerator, cut by only $0.4$--$9.3\%$).
Because the denominator falls faster than the numerator, PAR rises while
absolute peaks and cross-network imbalance improve--consistent with
distribution grids carrying distributed photovoltaics without stationary
storage.

\begin{table}[!t]
\footnotesize
\setlength{\tabcolsep}{4pt}
\caption{Per-Network Peak and PAR Decomposition (PV vs.\ PVB\_FULL).}
\label{tab:par}
\centering
\begin{tabular}{lccc}\toprule
Network & PV peak & FULL peak & PAR (PV$\to$FULL) \\ \midrule
Residential   & 650.6\,kW & 648.3\,kW & 1.152$\to$1.222 \\
Commercial    & 556.8\,kW & 545.3\,kW & 1.132$\to$1.176 \\
Industrial    & 1322.5\,kW & 1287.2\,kW & 1.456$\to$1.587 \\
Institutional & 787.8\,kW & 714.5\,kW & 1.654$\to$1.673 \\ \bottomrule
\end{tabular}
\end{table}

\paperfig[\columnwidth]{fig8_par_decomposition.png}%
 {Per-network peak total load (kW) under PV vs.\ PVB\_FULL, with the
 resulting peak-to-average ratio $D_{\max}/D_{\rm mean}$ in the lower
 panel, matching Table~\ref{tab:par}; PAR values use the canonical
 seed-42 evaluation.}%
 {fig:par}

\subsection{An Inherited Structural Limitation}\label{sec:inherited}
Across all four networks the PVB/PVB\_FULL EV-demand peak is reduced
relative to the uncoordinated EV baseline by a nearly identical margin
($\approx44.5\%$ with respect to the unbalanced EV peak;
Fig.~\ref{fig:uniformity}), despite the
networks' distinct diurnal shapes. Reproducing the base paper's own
two-network computation from its data and code gives the same uniformity:
the EV peak under PVB is a near-constant fraction ($\approx0.83$, within one
percentage point across the two networks) of the PV EV peak. The cause is
structural: independent per-network min-max normalization collapses
network-specific magnitude information before the demand equation is
applied, and the scale $\mu_i$ then re-scales each network by a fixed
fraction of its own unbalanced peak. The resulting EV peak reduction is
therefore approximately scale-invariant across networks. We document this as
a limitation \emph{inherited} from the base formulation rather than
introduced by our $N$-network, distance, or renewable extensions.

\paperfig[\columnwidth]{fig9_ev_peak_uniformity.png}%
{EV demand peak under PVB\_FULL as a fraction of the unbalanced EV peak for
each network--the near-uniform ($\approx$44--45\%) reduction inherited from
the base min--max normalization.}%
{fig:uniformity}

% ======================================================================
\section{Discussion}\label{sec:disc}

\subsection{Role and Scope of Distance Decay}
The measured contribution of distance decay ($-0.30\%$) quantifies spatial
friction at the specific scale tested ($1.16$--$7.06$~km). In such an
intra-city setting transit time is short, so the effect is small but
consistent and directionally correct. In larger metropolitan or multi-city
topologies, where travel time materially changes station choice
\cite{frade2011}, the term is expected to grow; the fitting of
$\alpha=0.15~\mathrm{km^{-1}}$ from revealed charging data is a natural
extension that we leave to future work.

\subsection{Inherited Normalization Limitation and Its Consequences}
The uniformity documented in Section~\ref{sec:inherited} is a genuine
property of the base demand equation, not an artifact of our
implementation. It means that the closed-form EV response is informative
about \emph{when} to move demand but weakly informative about \emph{which}
network carries the marginal kilowatt. The grid-utilization improvements
we report are nonetheless real, because the balancing signal itself
operates on raw utilizations whose scale information is preserved; only the
EV-demand *response* is scale-collapsed. A redesign that normalizes a
stacked multi-network vector jointly---or that preserves per-network
magnitudes through an affine, rather than min-max, transform---is the most
promising route to sharper cross-network shaping and is prioritized in
future work.

\subsection{Threats to Validity}
Residential and commercial loads and all meteorological inputs are real
Tetouan measurements; industrial and institutional profiles are synthetic,
constructed by re-weighting the real residential series. The statistical
test quantifies closed-form demand-noise sensitivity, not operational
stochasticity. The per-network PAR analysis inherits the assumption of no
stationary storage or export limits. Finally, the base paper's own data are
used for the inheritance analysis; we could not independently re-run its
DDPG training within this study's compute budget, so that portion of the
inheritance claim rests on re-computation with its reference implementation.

% ======================================================================
\section{Conclusion and Future Work}\label{sec:conc}
This paper extended a two-network collaborative DDPG EV charging pricing
framework along three directions: an exact, backward-compatible
generalization to arbitrary $N$ networks; a Haversine-distance exponential
decay governing cross-network driver willingness; and physics-based
solar/wind integration through a renewable-adjusted net load. On a
four-network Tetouan-grounded testbed, the combined framework reduced
cross-network utilization imbalance by $19.03\%$ relative to uncoordinated
pricing and by $15.28\%$ relative to the original pairwise formulation,
with the ranking preserved across four seasons, five inference-noise seeds,
and four DDPG training seeds. Ablation attributes most of the gain to
renewable net-load integration; generalization and distance awareness
contribute smaller, consistent improvements. We further documented, and
attributed to the base formulation, a structural EV peak-uniformity
property caused by independent min-max normalization.

Future work includes: (i) validation on additional real multi-feeder
telemetry; (ii) distance-decay studies on larger multi-city topologies
with calibrated $\alpha$; (iii) a redesigned demand-response equation that
preserves cross-network load-shape information; and (iv) extending the
statistical evaluation from inference-noise seeds to retrained DDPG
ensembles across a full annual cycle.

% ======================================================================
\section*{Acknowledgment}
The authors acknowledge the use of AI-assisted coding tools during the
implementation and verification stages of this work. All mathematical
formulations, experimental design decisions, and result interpretations
were made by the authors. The Tetouan smart-grid dataset was obtained from
the UCI Machine Learning Repository~\cite{salam2018}.

% ======================================================================
\appendix
\section{Numerical Equivalence Verification}\label{app:equiv}
To guard against relying on the algebraic argument of
Proposition~\ref{prop:n2} alone, the companion notebook verifies the
reduction numerically. The generalized function is called with $N=2$ and
its output is compared against the base pairwise function on the canonical
inputs.

\begin{algorithm}[H]
\caption{Numerical backward-compatibility check}
\begin{algorithmic}[1]
\State \textbf{input:} balancing functions $f_{\mathrm{pw}}$, $f_{\mathrm{gen}}$, load matrix $\mathbf{P}\in\mathbb{R}^{T\times2}$
\For{$t=1,\dots,T$}
  \State $s^{\mathrm{pw}}_1 \gets f_{\mathrm{pw}}(\util_1(t),\util_2(t))$
  \State $s^{\mathrm{gen}}_1 \gets f_{\mathrm{gen}}(\util_1(t),\util_2(t))$
\EndFor
\State \Return $\texttt{np.allclose}(s^{\mathrm{pw}}, s^{\mathrm{gen}},
\texttt{atol}=10^{-7})$
\end{algorithmic}
\end{algorithm}

%\section{Notebook Reproducibility}
%All tables are produced from the canonical notebook outputs in
%\texttt{data/} (files named \texttt{*\_canonical.csv}). Headline, ablation,
%and seasonal tables use the 5-seed inference protocol; the DDPG-robustness
%table uses the 4-seed training protocol.

% ======================================================================
\begin{thebibliography}{99}\IEEEoverridecommandlockouts%

\bibitem{lepolesa2026}
L. J. Lepolesa, K. E. Adetunji, K. Ouahada, Z. Liu, and L. Cheng,
``Dynamic electric vehicle charging pricing for load balancing in power
distribution networks based on collaborative DDPG agents,''
\emph{IEEE Trans. Smart Grid}, vol. 17, no. 4, pp. 3342--3353, Jul. 2026.

\bibitem{bailey2025}
M. R. Bailey, D. P. Brown, E. Myers, B. Shaffer, and F. A. Wolak,
``Electric vehicles and the energy transition: Unintended consequences of
time-of-use pricing,''
\emph{American Economic Review: Insights}, vol. 7, no. 4, pp. 550--566, 2025.

\bibitem{samende2025}
C. Samende, A. Akshaya, S. M. Bhagavathy, J. Melone, R. H. Snape, and
K. Stewart,
``Herding from uncoordinated smart charging of EVs: A real-life
demonstration,''
in \emph{Proc. IEEE PES ISGT Europe}, 2025.

\bibitem{jones2022}
C. B. Jones, W. Vining, M. Lave, T. Haines, C. Neuman, J. Bennett,
\emph{et al.},
``Impact of electric vehicle customer response to time-of-use rates on
distribution power grids,''
\emph{Energy Reports}, vol. 8, 2022.

\bibitem{faruqui2010}
A. Faruqui and S. Sergici,
``Household response to dynamic pricing of electricity: A survey of 15
experiments,''
\emph{J. Regulatory Economics}, vol. 38, no. 2, pp. 193--225, 2010.

\bibitem{borenstein2005}
S. Borenstein,
``The long-run efficiency of real-time electricity pricing,''
\emph{The Energy Journal}, vol. 26, no. 3, pp. 93--116, 2005.

\bibitem{clement2010}
K. Clement-Nyns, E. Haesen, and J. Driesen,
``The impact of charging plug-in hybrid electric vehicles on a residential
distribution grid,''
\emph{IEEE Trans. Power Syst.}, vol. 25, no. 1, pp. 371--380, 2010.

\bibitem{richardson2012tsg}
P. Richardson, D. Flynn, and A. Keane,
``Local versus centralized charging strategies for electric vehicles in low
voltage distribution systems,''
\emph{IEEE Trans. Smart Grid}, vol. 3, no. 2, pp. 1020--1028, 2012.

\bibitem{richardson2012tpwrs}
P. Richardson, D. Flynn, and A. Keane,
``Optimal charging of electric vehicles in low-voltage distribution
systems,''
\emph{IEEE Trans. Power Syst.}, vol. 27, no. 1, pp. 268--279, 2012.

\bibitem{oconnell2014}
A. O'Connell, D. Flynn, and A. Keane,
``Rolling multi-period optimization to control electric vehicle charging in
distribution networks,''
\emph{IEEE Trans. Power Syst.}, vol. 29, no. 1, pp. 340--348, 2014.

\bibitem{gan2013}
L. Gan, U. Topcu, and S. H. Low,
``Optimal decentralized protocol for electric vehicle charging,''
\emph{IEEE Trans. Power Syst.}, vol. 28, no. 2, pp. 940--951, 2013.

\bibitem{ma2013}
Z. Ma, D. S. Callaway, and I. A. Hiskens,
``Decentralized charging control of large populations of plug-in electric
vehicles,''
\emph{IEEE Trans. Control Syst. Technol.}, vol. 21, no. 1, pp. 67--78, 2013.

\bibitem{callaway2011}
D. S. Callaway and I. A. Hiskens,
``Achieving controllability of electric loads,''
\emph{Proc. IEEE}, vol. 99, no. 1, pp. 184--199, 2011.

\bibitem{baran1989}
M. E. Baran and F. F. Wu,
``Network reconfiguration in distribution systems for loss reduction and
load balancing,''
\emph{IEEE Trans. Power Del.}, vol. 4, no. 2, pp. 1401--1407, 1989.

\bibitem{hoke2013}
A. Hoke, R. Butler, J. Hambrick, and B. Kroposki,
``Steady-state analysis of maximum photovoltaic penetration levels on
typical distribution feeders,''
\emph{IEEE Trans. Sustain. Energy}, vol. 4, no. 2, pp. 350--357, 2013.

\bibitem{huff1964}
D. L. Huff,
``Defining and estimating a trading area,''
\emph{Journal of Marketing}, vol. 28, no. 3, pp. 34--38, 1964.

\bibitem{frade2011}
I. Frade, A. Ribeiro, G. Gon\c{c}alves, and A. P. Antunes,
``Optimal location of charging stations for electric vehicles in a
neighborhood in Lisbon, Portugal,''
\emph{Transp. Res. Rec.}, vol. 2252, no. 1, pp. 91--98, 2011.

\bibitem{hong2016}
T. Hong and S. Fan,
``Probabilistic electric load forecasting: A tutorial review,''
\emph{Int. J. Forecasting}, vol. 32, no. 3, pp. 914--938, 2016.

\bibitem{wang2016recency}
P. Wang, B. Liu, and T. Hong,
``Electric load forecasting with recency effect: A big data approach,''
\emph{Int. J. Forecasting}, vol. 32, no. 3, pp. 585--597, 2016.

\bibitem{salam2018}
A. Salam and A. El Hibaoui,
``Comparison of machine learning algorithms for the power consumption
prediction: Case study of Tetouan city,''
in \emph{Proc. 6th IEEE Int. Renewable Sustain. Energy Conf. (IRSEC)}, 2018;
dataset: \emph{Power Consumption of Tetouan City}, UCI Machine Learning
Repository, DOI \texttt{10.24432/C5B034}.

\bibitem{chen2016}
T. Chen and C. Guestrin,
``XGBoost: A scalable tree boosting system,''
in \emph{Proc. 22nd ACM SIGKDD Int. Conf. Knowledge Discovery and Data
Mining}, 2016, pp. 785--794.

\bibitem{ke2017}
G. Ke, Q. Meng, T. Finley, T. Wang, W. Chen, W. Ma, Q. Ye, and T.-Y. Liu,
``LightGBM: A highly efficient gradient boosting decision tree,''
in \emph{Advances in Neural Information Processing Systems}, 2017.

\bibitem{breiman2001}
L. Breiman,
``Random forests,''
\emph{Machine Learning}, vol. 45, no. 1, pp. 5--32, 2001.

\bibitem{hochreiter1997}
S. Hochreiter and J. Schmidhuber,
``Long short-term memory,''
\emph{Neural Computation}, vol. 9, no. 8, pp. 1735--1780, 1997.

\bibitem{lillicrap2016}
T. P. Lillicrap \emph{et al.},
``Continuous control with deep reinforcement learning,''
in \emph{Proc. Int. Conf. Learning Representations (ICLR)}, 2016.

\bibitem{mnih2015}
V. Mnih \emph{et al.},
``Human-level control through deep reinforcement learning,''
\emph{Nature}, vol. 518, pp. 529--533, 2015.

\bibitem{haarnoja2018}
T. Haarnoja, A. Zhou, P. Abbeel, and S. Levine,
``Soft actor-critic: Off-policy maximum entropy deep reinforcement learning
with a stochastic actor,''
in \emph{Proc. 35th Int. Conf. Machine Learning (ICML)}, 2018,
pp. 1865--1874.

\bibitem{schulman2017}
J. Schulman, F. Wolski, P. Dhariwal, A. Radford, and O. Klimov,
``Proximal policy optimization algorithms,''
arXiv preprint arXiv:1707.06347, 2017.

\bibitem{kingma2015}
D. P. Kingma and J. Ba,
``Adam: A method for stochastic optimization,''
in \emph{Proc. Int. Conf. Learning Representations (ICLR)}, 2015.

\bibitem{vazquez2019}
J. R. V\u{a}zquez-Canteli and Z. Nagy,
``Reinforcement learning for demand response: A review of algorithms and
modeling techniques,''
\emph{Applied Energy}, vol. 235, pp. 1072--1089, 2019.

\bibitem{qiu2020tia}
D. Qiu, Y. Ye, D. Papadaskalopoulos, and G. Strbac,
``A deep reinforcement learning method for pricing electric vehicles with
discrete charging levels,''
\emph{IEEE Trans. Ind. Appl.}, vol. 56, no. 5, pp. 5901--5912, 2020.

\bibitem{liu2021access}
D. Liu, W. Wang, L. Wang, H. Jia, and M. Shi,
``Dynamic pricing strategy of electric vehicle aggregators based on DDPG
reinforcement learning algorithm,''
\emph{IEEE Access}, vol. 9, pp. 21556--21566, 2021.

\bibitem{zhao2022tte}
Z. Zhao and C. K. M. Lee,
``Dynamic pricing for EV charging stations: A deep reinforcement learning
approach,''
\emph{IEEE Trans. Transp. Electrific.}, vol. 8, no. 2, pp. 2456--2468, 2021.

\bibitem{ye2022}
Y. Ye, D. Papadaskalopoulos, Q. Yuan, Y. Tang, and G. Strbac,
``Multi-agent deep reinforcement learning for coordinated energy trading and
flexibility services provision in local electricity markets,''
\emph{IEEE Trans. Smart Grid}, vol. 14, no. 2, pp. 1541--1554, 2022.

\bibitem{yang2024}
X. Yang, T. Cui, H. Wang, and Y. Ye,
``Multiagent deep reinforcement learning for electric vehicle fast charging
station pricing game in electricity-transportation nexus,''
\emph{IEEE Trans. Ind. Informat.}, vol. 20, no. 4, pp. 6345--6355, 2024,
doi: 10.1109/TII.2023.3345457.

\bibitem{lai2022tsg}
S. Lai, J. Qiu, Y. Tao, and J. Zhao,
``Pricing for electric vehicle charging stations based on the responsiveness
of demand,''
\emph{IEEE Trans. Smart Grid}, vol. 14, no. 1, pp. 530--544, 2023.

\bibitem{cui2021tsg}
Y. Cui, Z. Hu, and X. Duan,
``Optimal pricing of public electric vehicle charging stations considering
operations of coupled transportation and power systems,''
\emph{IEEE Trans. Smart Grid}, vol. 12, no. 4, pp. 3278--3288, 2021.

\bibitem{aljafari2023}
B. Aljafari, P. R. Jeyaraj, A. C. Kathiresan, and S. B. Thanikanti,
``Electric vehicle optimum charging-discharging scheduling with dynamic
pricing employing multi agent deep neural network,''
\emph{Computers and Electrical Engineering}, vol. 105, Art. no. 108555, 2023.

\bibitem{yan2021tsg}
L. Yan, X. Chen, J. Zhou, Y. Chen, and J. Wen,
``Deep reinforcement learning for continuous electric vehicles charging
control with dynamic user behaviors,''
\emph{IEEE Trans. Smart Grid}, vol. 12, no. 6, pp. 5124--5134, 2021.

\bibitem{paudel2023}
D. Paudel and T. K. Das,
``A deep reinforcement learning approach for power management of
battery-assisted fast-charging EV hubs participating in day-ahead and
real-time electricity markets,''
\emph{Energy}, vol. 283, Art. no. 129097, 2023.

\end{thebibliography}

\end{document}